\section{The multiset hierarchy}

We first record a general majorization device.

\begin{lemma}[Multiset majorant]\label{lem:multmaj}
Fix $r\geq1$ and suppose $c,\alpha,\beta>0$ satisfy, for every integer $m\geq1$,
\begin{align}
k_c(m)&\leq \alpha m+\beta m\ind_{m\leq r},\label{eq:onmaj}\\
c^2&\leq 2\alpha m+2\beta m\ind_{m\leq r}.\label{eq:pairmaj}
\end{align}
Then
\begin{equation}\label{eq:multgeneric}
\liminf_{T\to\infty}\frac{M_{\leq r}(T,2T)}{N(T,2T)}
\geq \frac{2c-\kMT-\alpha}{\beta}.
\end{equation}
\end{lemma}

\begin{proof}
Sum \eqref{eq:onmaj} over critical-line points and \eqref{eq:pairmaj} over reflected pairs.  By \eqref{eq:countidentity}, the $\alpha$-terms sum exactly to $\alpha N$ and the $\beta$-terms to $\beta M_{\leq r}$.  Proposition~\ref{prop:master} gives, for every $\varepsilon>0$ and all sufficiently large $T$,
\[
(2c-\kMT-\varepsilon)N\leq \alpha N+\beta M_{\leq r}.
\]
Divide by $N$, take the lower limit, and then let $\varepsilon\downarrow0$.
\end{proof}

\begin{proof}[Proof of Theorem~\ref{thm:multiset}]
For $r=1$, take
\[
c=2,\qquad \alpha=2,\qquad \beta=1.
\]
Then $k_2(1)=3=2+1$, while $k_2(m)=4\leq2m$ for $m\geq2$; the pair inequality is immediate.  Lemma~\ref{lem:multmaj} gives $2-\kMT$.

For $r=2$, take
\[
c=3,\qquad \alpha=3,\qquad \beta=2.
\]
Here $k_3(1)=5$, $k_3(2)=8$, and $k_3(m)=9$ for $m\geq3$, so the two pointwise inequalities again hold.  This gives $(3-\kMT)/2$.

Now let $r\geq3$ and set
\begin{equation}\label{eq:Cdef}
C:=2+\sqrt2,
\qquad
\alpha_r:=\frac{C^2}{r+1},
\qquad
\beta_r:=\frac{C^2(r-1)}{2(r+1)}.
\end{equation}
The balancing identity
\begin{equation}\label{eq:balance}
C^2=2(2C-1)
\end{equation}
implies
\[
\alpha_r+\beta_r=2C-1=\frac{C^2}{2}.
\]
For $m\leq r$ and $m\leq C$, equation \eqref{eq:kc} gives
\[
k_C(m)=m(2C-m)\leq m(2C-1)=m(\alpha_r+\beta_r).
\]
For $m\leq r$ and $m\geq C$, one has
\[
k_C(m)=C^2\leq \frac{C^2}{2}m=m(\alpha_r+\beta_r).
\]
For $m\geq r+1$, since $r+1\geq4>C$,
\[
k_C(m)=C^2=\alpha_r(r+1)\leq\alpha_rm.
\]
This proves the on-line majorant.  For a reflected pair, if $m\leq r$ then
\[
2m(\alpha_r+\beta_r)=mC^2\geq C^2,
\]
and if $m\geq r+1$ then $2\alpha_rm\geq2C^2>C^2$.  Hence Lemma~\ref{lem:multmaj} applies and yields
\[
\liminf_{T\to\infty}\frac{M_{\leq r}(T,2T)}{N(T,2T)}
\geq
\frac{2C-\kMT-\alpha_r}{\beta_r}.
\]
Direct simplification, using $(2+\sqrt2)^{-2}=3/2-\sqrt2$, gives
\[
\frac{2C-\kMT-\alpha_r}{\beta_r}
=1-(3-2\sqrt2)(\kMT-1)\frac{r+1}{r-1},
\]
which is \eqref{eq:Mr}.

The cumulative result follows by dyadic summation exactly as in \cite{Claude2026}.  The fixed primitive Dirichlet $L$-function case uses the corresponding trace theorem from that paper; the zero-side scalar argument is unchanged.
\end{proof}
